\chapter{ارزیابی و نتایج}
\sisetup{round-mode=places, round-precision=3}
\section{مقدمه}

این فصل به ارائه نتایج حاصل از اجرای مدل‌های پیش‌بینی طراحی‌شده برای تخمین اولین رکورد شیردهی گاوهای شیری اختصاص دارد. همان‌طور که در فصل‌های پیشین اشاره شد، هدف اصلی این پژوهش توسعه الگوریتم‌هایی مبتنی بر یادگیری ماشین برای تحلیل داده‌های تاریخی و پیش‌بینی عملکرد شیر در دوره انتقال است. در این راستا، مجموعه‌ای از مدل‌ها شامل رگرسیون ریج، جنگل‌های تصادفی، رگرسیون بردار پشتیبانی، افزایش گرادیان افراطی، تقویت دسته‌ای، و شبکه‌های عصبی چندلایه و عمیق مورد ارزیابی قرار گرفته‌اند. این مدل‌ها با استفاده از داده‌های جمع‌آوری‌شده از 345676 رأس گاو در 99 گله طی سال‌های 1390 تا 1401، که به سه بخش آموزشی، آزمایشی، و اعتبارسنجی تقسیم شده‌اند، آموزش داده شده‌اند.

در این فصل، ابتدا نتایج کمی حاصل از اجرای مدل‌ها با استفاده از شش معیار ارزیابی شامل ضریب تعیین، میانگین قدرمطلق خطا، ریشه میانگین مربع خطا، درصد خطای میانگین، میانگین قدرمطلق درصد خطای متقارن، و انحراف استاندارد نسبی ارائه می‌شود. سپس، با تحلیل این نتایج، مدل بهینه انتخاب شده و به‌صورت تفصیلی مورد بررسی قرار می‌گیرد. هدف از این تحلیل، نه‌تنها شناسایی بهترین عملکرد پیش‌بینی، بلکه ارائه بینش‌هایی برای بهبود مدیریت دوره انتقال و تصمیم‌گیری‌های مبتنی بر داده در صنعت دامداری است. در نهایت، نتایج این فصل به‌عنوان پایه‌ای برای بحث‌های آینده در پژوهش‌های آتی، به‌ویژه در توسعه و بهینه‌سازی مدل‌ها، مورد استفاده قرار خواهد گرفت.

\section{تحلیل داده‌ها و نتایج مدل‌ها}

\begin{enumerate}
	\item \textbf{ضریب تعیین\LTRfootnote{$R^2$}}: نشان‌دهنده درصد واریانس توضیح‌داده‌شده توسط مدل است و مقداری بین 0 تا 1 دارد. مقادیر بالاتر نشان‌دهنده برازش بهتر مدل هستند.
	
	\item \textbf{میانگین قدرمطلق خطا\LTRfootnote{MAE}}: میانگین اختلاف مطلق بین مقادیر پیش‌بینی‌شده و واقعی را نشان می‌دهد و واحد آن کیلوگرم است.
	
	\item \textbf{ریشه میانگین مربع خطا\LTRfootnote{RMSE}}: حساس به خطاهای بزرگ است و نشان‌دهنده پراکندگی خطاها به‌صورت ریشه دوم است.
	
	\item \textbf{درصد خطای میانگین\LTRfootnote{MPE}}: جهت کلی خطاها را نشان می‌دهد و به‌صورت درصدی بیان می‌شود.
	
	\item \textbf{میانگین قدرمطلق درصد خطای متقارن\LTRfootnote{sMAPE}}: برای مقایسه نسبی خطاها در مقیاس درصد استفاده می‌شود.
	
	\item \textbf{انحراف استاندارد نسبی\LTRfootnote{SDR}}: نسبت ریشه میانگین مربع خطا به انحراف معیار مشاهدات واقعی است و نشان‌دهنده میزان پراکندگی خطا نسبت به داده‌ها است.
\end{enumerate}

نتایج حاصل از اجرای مدل‌ها در جدول زیر ارائه شده است:

\begin{table}[H]
	\centering
	\caption{مقایسه عملکرد مدل‌ها بر اساس معیارهای ارزیابی}
	\begin{LTR}
		\begin{adjustbox}{width=\columnwidth}
			\rowcolors{2}{customblue!20}{white} % رنگ جدید پیشنهادی
			\begin{tabular}{
					>{\centering\arraybackslash}m{4.3cm} % ستون اول با وسط‌چینی عمودی
					*{6}{>{\centering\arraybackslash}m{2.2cm}} % سایر ستون‌ها
				}
				\toprule
				\rowcolor{customblue}
				\textbf{مدل} & \RLE{\textbf{ضریب تعیین}} & \RLE{\textbf{میانگین قدرمطلق خطا}} & \RLE{\textbf{ریشه میانگین مربع خطا}} & \RLE{\textbf{درصد خطای میانگین}} & \RLE{\textbf{میانگین قدرمطلق درصد خطای متقارن}} & \RLE{\textbf{انحراف استاندارد نسبی}} \\
				\midrule
				Ridge\_Initial  & \num{0.310994563} & \num{7.027555958} & \num{9.206984387} & -\num{8.085322509} & \num{17.68397569} & \num{0.516879496} \\
				Ridge\_Tuned  & \num{0.313394604} & \num{7.00567259} & \num{9.190934868} & -\num{7.877926746} & \num{17.66417559} & \num{0.557122089} \\
				Random Forest\_Initial  & \num{0.347708828} & \num{6.828522689} & \num{8.958325384} & -\num{6.93634364} & \num{17.22296862} & \num{0.610603574} \\
				Random Forest\_Tuned  & \num{0.357528448} & \num{6.759105101} & \num{8.890640168} & -\num{7.31582344} & \num{17.05506122} & \num{0.60316898} \\
				Random Forest with Final Cross-Validation & \num{0.352553746} & \num{6.784989338} & \num{8.924994218} & -\num{7.61062843} & \num{17.09271223} & \num{0.574371254} \\
				Random Forest\_Initial\_with\_grp  & \num{0.340053875} & \num{6.878538826} & \num{9.01073721} & -\num{6.964281014} & \num{17.34244231} & \num{0.613633829} \\
				Random Forest\_Tuned\_with\_grp  & \num{0.351373163} & \num{6.792522921} & \num{8.933127636} & -\num{7.608877095} & \num{17.11226016} & \num{0.574870127} \\
				Random Forest\_with Final Cross-Validation\_grp & \num{0.351373163} & \num{6.792522921} & \num{8.933127636} & -\num{7.608877095} & \num{17.11226016} & \num{0.574870127} \\
				SVR\_Linear  & \num{0.307137444} & \num{6.970674056} & \num{9.23271922} & -\num{9.731247682} & \num{17.58411057} & \num{0.584671629} \\
				SVR\_Polynomial  & \num{0.264796863} & \num{7.216205775} & \num{9.510640332} & -\num{9.820527027} & \num{18.17611214} & \num{0.554772935} \\
				SVR\_RBF  & \num{0.364292439} & \num{6.637657728} & \num{8.843715657} & -\num{9.237820776} & \num{16.75859351} & \num{0.615623665} \\
				SVR\_RBF\_Robust  & \num{0.360090177} & \num{6.643834386} & \num{8.86102783} & -\num{9.240319703} & \num{16.77569945} & \num{0.616597688} \\
				XGBoost  & \num{0.366754141} & \num{6.691355247} & \num{8.82657593} & -\num{7.290271157} & \num{16.90602028} & \num{0.617812527} \\
				XGBoost\_Tuned  & \num{0.372973212} & \num{6.660094991} & \num{8.783126339} & -\num{7.318830137} & \num{16.82175142} & \num{0.60636653} \\
				CatBoost\_Initial  & \num{0.401540057} & \num{6.438195384} & \num{8.504328151} & -\num{6.66673337} & \num{16.21511352} & \num{0.628887986} \\
				CatBoost\_Tuned  & \num{0.38592024} & \num{6.538945512} & \num{8.614594868} & -\num{6.857419311} & \num{16.45040884} & \num{0.607961385} \\
				CatBoost\_Initial\_Standard  & \num{0.372655433} & \num{6.658741634} & \num{8.785351717} & -\num{7.343736893} & \num{16.81562966} & \num{0.604141488} \\
				MLP & \num{0.370086641} & \num{6.654381139} & \num{8.803320059} & -\num{7.64569314} & \num{16.80819888} & \num{0.620456942} \\
				DNN & \num{0.36876709} & \num{6.692697966} & \num{8.812535893} & -\num{6.888868896} & \num{16.89858016} & \num{0.607973874} \\
				\bottomrule
			\end{tabular}
		\end{adjustbox}
	\end{LTR}
\end{table}

بر اساس جدول فوق، مدل‌های مختلف عملکرد متفاوتی از خود نشان داده‌اند. به‌طور کلی، مدل تقویت دسته‌ای اولیه با ضریب تعیین \num{0.4015} بالاترین برازش را داشته و میانگین قدرمطلق خطای آن (\num{6.4382} کیلوگرم) کمترین مقدار را در بین مدل‌ها داراست. این نشان‌دهنده توانایی این مدل در پیش‌بینی دقیق‌تر مقادیر واقعی است. از سوی دیگر، مدل رگرسیون بردار پشتیبانی با هسته چندجمله‌ای با ضریب تعیین \num{0.2648} کمترین عملکرد را داشته و خطاهای بالاتری (ریشه میانگین مربع خطا \num{9.5106} و میانگین قدرمطلق خطا \num{7.2162}) را نشان می‌دهد، که می‌تواند ناشی از پیچیدگی بیش‌ازحد یا ناسازگاری هسته با داده‌های موجود باشد.

\section{مقایسه و انتخاب مدل بهینه}

برای انتخاب مدل بهینه، معیارهای مختلف به‌صورت ترکیبی بررسی شده‌اند. ضریب تعیین به‌عنوان شاخص اصلی برازش، میانگین قدرمطلق خطا به‌عنوان معیاری از دقت پیش‌بینی، و ریشه میانگین مربع خطا به‌عنوان حساسیت به خطاهای بزرگ، در اولویت قرار گرفته‌اند. همچنین، انحراف استاندارد نسبی، درصد خطای میانگین، و میانگین قدرمطلق درصد خطای متقارن برای ارزیابی تعادل و ثبات مدل‌ها مورد توجه بوده‌اند.

برای نشان دادن برتری مدل بهینه (تقویت دسته‌ای اولیه)، جدول زیر درصد بهبود یا اختلاف این مدل نسبت به سایر مدل‌ها بر اساس تمامی معیارهای ارزیابی ارائه شده است. مقادیر مثبت در ضریب تعیین نشان‌دهنده برازش بهتر و در معیارهای خطا نشان‌دهنده کاهش خطا یا بهبود عملکرد هستند.
\begin{table}[H]
	\centering
	\caption{درصد بهبود مدل بهینه (تقویت دسته‌ای اولیه) نسبت به سایر مدل‌ها}
	\begin{LTR}
		\begin{adjustbox}{width=\columnwidth}
			\rowcolors{2}{customblue!20}{white}
			\begin{tabular}{
					>{\centering\arraybackslash}m{4.3cm}
					*{6}{>{\centering\arraybackslash}m{2cm}}
				}
				\toprule
				\rowcolor{customblue}
				\textbf{مدل} & \RLE{\textbf{بهبود ضریب تعیین (\%)}} & \RLE{\textbf{کاهش MAE (\%)}} & \RLE{\textbf{کاهش RMSE (\%)}} & \RLE{\textbf{کاهش MPE (\%)}} & \RLE{\textbf{کاهش sMAPE (\%)}} & \RLE{\textbf{کاهش SDR (\%)}} \\
				\midrule
				Ridge\_Initial  & \num{29.16} & \num{8.38} & \num{7.65} & \num{17.52} & \num{8.33} & - \num{21.67} \\
				Ridge\_Tuned  & \num{28.14} & \num{8.12} & \num{7.47} & \num{15.37} & \num{8.26} & - \num{0.11} \\
				Random Forest\_Initial  & \num{15.47} & \num{5.71} & \num{5.05} & \num{3.89} & \num{5.87} &  - \num{2.98} \\
				Random Forest\_Tuned  & \num{12.31} & \num{4.74} & \num{4.32} & \num{8.86} & \num{5.14} & - \num{0.42} \\
				Random Forest with Final Cross-Validation & \num{13.91} & \num{5.11} & \num{4.68} & \num{12.40} & \num{5.10} & - \num{9.47} \\
				Random Forest\_Initial\_with\_grp  & \num{18.12} & \num{6.40} & \num{5.61} & \num{4.23} & \num{6.49} & - \num{1.25} \\
				Random Forest\_Tuned\_with\_grp  & \num{14.31} & \num{5.23} & \num{4.75} & \num{12.37} & \num{5.25} & - \num{9.59} \\
				Random Forest\_with Final Cross-Validation\_grp & \num{14.31} & \num{5.23} & \num{4.75} & \num{12.37} & \num{5.25} & - \num{9.59} \\
				SVR\_Linear  & \num{30.77} & \num{7.66} & \num{7.90} & \num{31.49} & \num{7.79} & - \num{7.52} \\
				SVR\_Polynomial  & \num{51.71} & \num{10.80} & \num{10.58} & \num{32.12} & \num{10.80} & - \num{13.34} \\
				SVR\_RBF  & \num{10.28} & \num{3.02} & \num{3.83} & \num{27.85} & \num{3.24} & \num{2.14} \\
				SVR\_RBF\_Robust  & \num{11.61} & \num{3.10} & \num{4.00} & \num{27.89} & \num{3.32} & \num{1.95} \\
				XGBoost  & \num{9.50} & \num{3.77} & \num{3.62} & \num{8.51} & \num{4.10} & \num{1.76} \\
				XGBoost\_Tuned  & \num{7.68} & \num{3.33} & \num{3.14} & \num{8.89} & \num{3.63} & - \num{3.67} \\
				CatBoost\_Tuned  & \num{4.12} & \num{1.54} & \num{1.28} & \num{2.77} & \num{1.43} & - \num{3.42} \\
				CatBoost\_Initial\_Standard  & \num{7.79} & \num{3.32} & \num{3.19} & \num{9.24} & \num{3.54} & - \num{4.02} \\
				MLP  & \num{8.47} & \num{3.28} & \num{3.34} & \num{12.82} & \num{3.52} & \num{1.28} \\
				DNN  & \num{8.91} & \num{3.81} & \num{3.47} & \num{3.23} & \num{4.07} & - \num{3.48} \\
				\bottomrule
			\end{tabular}
		\end{adjustbox}
	\end{LTR}
\end{table}

بررسی جدول نشان می‌دهد که مدل تقویت دسته‌ای اولیه با بالاترین ضریب تعیین (\num{0.4015}) و کمترین میانگین قدرمطلق خطا (\num{6.4382}) بهترین عملکرد را دارد. این مدل نسبت به سایر مدل‌ها، از جمله نسخه‌های بهینه‌شده خود (مانند تقویت دسته‌ای بهینه‌شده با ضریب تعیین \num{0.3859})، ثبات و دقت بیشتری نشان می‌دهد. همچنین، ریشه میانگین مربع خطا این مدل (\num{8.5043}) نسبت به میانگین ریشه میانگین مربع خطا سایر مدل‌ها (حدود \num{8.9}) پایین‌تر است، که نشان‌دهنده کاهش خطاهای بزرگ در پیش‌بینی‌هاست. با این حال، انحراف استاندارد نسبی این مدل (\num{0.6289}) کمی بالاتر از برخی مدل‌های دیگر (مانند رگرسیون ریج اولیه با \num{0.5169}) است، که می‌تواند به پراکندگی بیشتر داده‌ها یا حساسیت مدل به نویز نسبت داده شود.

مدل‌های شبکه عصبی (شبکه عصبی چندلایه و عمیق) با ضریب تعیین حدود \num{0.37} و میانگین قدرمطلق خطا بین \num{6.65} و \num{6.69} عملکرد متوسطی داشته‌اند، اما به دلیل پیچیدگی بالاتر و نیاز به تنظیم دقیق‌تر پارامترها، در مقایسه با تقویت دسته‌ای اولیه کمتر بهینه به نظر می‌رسند. از سوی دیگر، مدل‌های رگرسیون بردار پشتیبانی، به‌ویژه با هسته‌های غیرخطی (چندجمله‌ای و شعاعی)، به دلیل ضریب تعیین پایین و خطاهای بالاتر، برای این داده‌ها مناسب تشخیص داده نشده‌اند.

بنابراین، با توجه به معیارهای فوق و تعادل بین دقت، برازش، و سادگی، مدل تقویت دسته‌ای اولیه به‌عنوان مدل بهینه انتخاب می‌شود. این انتخاب با در نظر گرفتن کاربرد عملی مدل در صنعت دامداری، که نیاز به دقت بالا و تعمیم‌پذیری مناسب دارد، انجام شده است.

\section{تحلیل تفصیلی مدل بهینه}

مدل تقویت دسته‌ای اولیه، که بر پایه الگوریتم تقویت گرادیانی\LTRfootnote{Gradient Boosting} طراحی شده است، با استفاده از داده‌های آموزشی و پارامترهای پیش‌فرض خود اجرا شده است. این مدل به‌صورت خودکار متغیرهای دسته‌ای را پردازش کرده و با بهره‌گیری از روش آمار هدف، نیازی به کدگذاری دستی ندارد. ساختار این مدل شامل مجموعه‌ای از درخت‌های تصمیم است که به‌صورت ترتیبی ساخته شده و هر درخت خطاهای مدل قبلی را اصلاح می‌کند.

\subsection{عملکرد کمی مدل}
بر اساس نتایج جدول، این مدل ضریب تعیین \num{0.4015} را به دست آورده است، به این معنا که حدود 40 درصد از واریانس داده‌های واقعی توسط این مدل توضیح داده شده است. این مقدار، اگرچه بهینه‌ترین مقدار ممکن نیست، اما در مقایسه با سایر مدل‌ها نشان‌دهنده برازش قابل قبولی است. میانگین قدرمطلق خطا (\num{6.4382} کیلوگرم) نشان می‌دهد که پیش‌بینی‌های این مدل به‌طور متوسط کمتر از \num{6.5} کیلوگرم با مقادیر واقعی اختلاف دارند، که برای پیش‌بینی تولید روزانه شیر در بازه‌ای که معمولاً بین 20 تا 40 کیلوگرم متغیر است، قابل قبول است. ریشه میانگین مربع خطا (\num{8.5043}) نیز تأییدکننده این است که خطاهای بزرگ در پیش‌بینی‌ها به حداقل رسیده‌اند.

درصد خطای میانگین (\num{6.6667}- درصد) نشان‌دهنده تمایل مدل به پیش‌بینی مقادیری کمی بالاتر از مقادیر واقعی است، که می‌تواند ناشی از نویز موجود در داده‌ها یا ویژگی‌های نادیده‌گرفته‌شده باشد. میانگین قدرمطلق درصد خطای متقارن (\num{16.2151} درصد) نیز بیانگر ثبات نسبی مدل در پیش‌بینی‌های نسبی است. انحراف استاندارد نسبی (\num{0.6289}) نشان می‌دهد که پراکندگی خطاها نسبت به انحراف معیار داده‌ها در سطح متوسطی قرار دارد.

\subsection{تحلیل کیفی و ویژگی‌های مؤثر}
برای درک بهتر عملکرد مدل، اهمیت ویژگی‌ها در پیش‌بینی‌ها بررسی شده است. داده‌های تاریخی تولید شیر، به‌ویژه تولید 305 روزه شیردهی قبلی و روزهای شیردهی در اولین تست، بیشترین تأثیر را بر نتایج داشته‌اند. این یافته با مطالعات پیشین (مانند \cite{Salamone2022}) هم‌راستا است و نشان می‌دهد که سابقه تولید دام نقش کلیدی در پیش‌بینی عملکرد آتی دارد. همچنین، متغیرهایی مانند فصل زایش و میانگین عملکرد گله به‌عنوان عوامل محیطی، در کاهش خطاها مؤثر بوده‌اند.

\subsection{محدودیت‌ها و نقاط قوت مدل}
نقطه قوت اصلی این مدل، سادگی نسبی و توانایی آن در پردازش داده‌های دسته‌ای بدون نیاز به پیش‌پردازش پیچیده است. با این حال، انحراف استاندارد نسبی بالاتر نسبت به برخی مدل‌ها (مانند رگرسیون ریج) می‌تواند نشان‌دهنده حساسیت مدل به نویز یا نیاز به تنظیم بیشتر پارامترها باشد. این موضوع در پژوهش‌های آینده باید مورد بررسی قرار بگیرد.

\section{چالش‌ها و محدودیت‌ها}

این پژوهش به دلیل ماهیت صنعتی و غیرآکادمیک خود با چالش‌های متعددی همراه بود. یکی از اصلی‌ترین مشکلات، فرآیند جمع‌آوری دیتاست بود که نیازمند هماهنگی‌های گسترده، برگزاری جلسات متعدد و پیگیری‌های مستمر بود. دیتاست شامل اطلاعات حساس و مهم از دامداری‌های سراسر کشور بود، که این موضوع همکاری با نهادهای مربوطه را پیچیده می‌کرد. علاوه بر این، فهم دقیق مسئله نیز چالش‌برانگیز بود؛ به‌عنوان دانشجویان مهندسی کامپیوتر، هیچ‌گونه آشنایی اولیه با مفاهیم دامداری، اصطلاحات مرتبط یا حوزه‌های تخصصی این صنعت نداشتیم. حتی در جلسات توضیحی که برگزار می‌شد، مواجهه با اصطلاحات ناآشنا و عدم درک کامل موضوع، فرآیند درک مسئله را دشوارتر می‌کرد.

یکی دیگر از موانع، دسترسی به داده‌ها بود. نهادهای مسئول به دلیل نگرانی از افشای اطلاعات یا تصور اینکه ممکن است جای آن‌ها را بگیریم، با اکراه اطلاعات را در اختیار ما قرار می‌دادند، که این امر زمان و انرژی زیادی را صرف مذاکره و جلب اعتماد کرد. در اواسط پژوهش نیز با مشکلی اساسی مواجه شدیم؛ متوجه شدیم که ویژگی "شیر تجمعی 305 روز" که در دیتاست ارائه‌شده بود، خود از متغیر هدف مشتق شده و نباید در مدل استفاده می‌شد. این کشف منجر به ابطال تلاش‌های قبلی برای آموزش مدل‌ها شد و ما را مجبور کرد کل فرآیند را از ابتدا بازطراحی و اجرا کنیم.

کیفیت دیتاست نیز از دیگر محدودیت‌های این پژوهش بود. بر اساس مطالعات مشابه که در مراجع ذکر شده‌اند، ویژگی‌های اضافی (مانند داده‌های سلامتی یا شرایط محیطی خاص) معمولاً برای پیش‌بینی شیر استفاده می‌شوند، اما این ویژگی‌ها در دیتاست ما موجود نبودند. همچنین، دیتاست از استانداردهای یک دیتاست آکادمیک تر و تمیز برخوردار نبود و شامل نویز و ناقص‌بودن‌هایی بود که تحلیل را پیچیده‌تر کرد. با وجود این چالش‌ها، نتایج به‌دست‌آمده نشان‌دهنده پتانسیل مدل پیشنهادی در شرایط واقعی است، اما برای بهبود بیشتر، نیاز به دیتای جامع‌تر و باکیفیت‌تر احساس می‌شود.


\section{نتیجه‌گیری}

در این فصل، نتایج حاصل از اجرای 19 مدل مختلف برای پیش‌بینی اولین رکورد شیردهی گاوهای شیری ارائه شد. تحلیل‌ها نشان داد که مدل تقویت دسته‌ای اولیه با ضریب تعیین \num{0.4015} و میانگین قدرمطلق خطا \num{6.4382} کیلوگرم، بهترین عملکرد را در بین مدل‌ها دارد و به‌عنوان مدل بهینه انتخاب شده است. این مدل با استفاده از داده‌های تاریخی و متغیرهای کلیدی، توانسته است تعادلی مناسب بین دقت و تعمیم‌پذیری ایجاد کند. تحلیل تفصیلی این مدل نشان داد که داده‌های تولید قبلی و شرایط محیطی نقش مهمی در پیش‌بینی‌ها ایفا می‌کنند.

نتایج این فصل، نه‌تنها پایه‌ای برای ارزیابی مدیریت دوره انتقال فراهم می‌کند، بلکه مسیر را برای توسعه بیشتر مدل‌ها در فصل‌های بعدی هموار می‌سازد. مدل بهینه نسبت به مدل آماری فعلی مورد استفاده در تعاونی دامداران استان اصفهان، که ضریب تعیین آن حدود \num{0.30} است، حدود 10 درصد بهبود در برازش (ضریب تعیین) نشان می‌دهد. این برتری، ناشی از توانایی مدل پیشنهادی در پردازش داده‌های پیچیده و استخراج الگوهای پنهان از داده‌های تاریخی است، در حالی که مدل آماری فعلی به دلیل محدودیت‌های روش‌های کلاسیک، دقت کمتری در پیش‌بینی‌های دقیق دارد. در ادامه، با تمرکز بر بهینه‌سازی پارامترها و گنجاندن داده‌های سلامتی، انتظار می‌رود دقت پیش‌بینی‌ها بهبود یابد. این یافته‌ها می‌توانند به دامداران در اتخاذ تصمیمات بهینه برای افزایش بهره‌وری و سلامت گله کمک کنند.
